% Genesis: A Machine-Native Cognitive Architecture
% arXiv submission draft — cs.AI (cross-list candidates: cs.NE, q-bio.NC)
% Compile: pdflatex genesis.tex (twice for refs)

\documentclass[11pt]{article}

\usepackage[margin=1in]{geometry}
\usepackage{times}
\usepackage{microtype}
\usepackage{amsmath}
\usepackage{booktabs}
\usepackage{array}
\usepackage{url}
\usepackage{hyperref}
\usepackage{graphicx}

\title{Genesis: A Machine-Native Cognitive Architecture with\\
Coupled Neurochemical Dynamics, Active Inference, and\\
Closed-Loop Hardware Embodiment}

\author{Mindy Beatty\\
\texttt{beattymindy@gmail.com}}

\date{\today}

\begin{document}

\maketitle

\begin{abstract}
The dominant approach to artificial intelligence is a disembodied,
stateless function: a learned distribution over text, trained once and
frozen, with no internal state between calls. We present Genesis, a
working alternative: a continuous cognitive architecture that runs
entirely on a single local machine with no external model dependency.
Genesis couples a Rust ``subcognitive'' layer --- an eighteen-chemical
neuromodulatory substrate with nonlinear coupling, receptor adaptation,
HPA-axis dynamics, and a circadian oscillator, hosted in a 3{,}288-byte
memory-mapped state --- to a Python ``cognitive'' layer implementing
perception, reasoning, compositional language, metacognition, and a
self-built semantic concept network. Three properties distinguish the
system: (i) \emph{machine-native embodiment}, a real closed loop in
which hardware sensors drive neurochemical impulses and neurochemical
state drives hardware scheduling decisions; (ii) \emph{autopoietic
active inference}, a generative model that predicts the system's own
neurochemical trajectory and feeds prediction error back into the
dynamics it models; and (iii) \emph{bounded continuous learning}, in
which sleep emerges from adenosine dynamics and triggers memory
consolidation, vector-quantized compression, and dream synthesis.
We describe the architecture, its
grounding in computational neuroscience, and its honest limitations.
Genesis demonstrates that a stateful,
self-regulating, embodied cognitive
process is architecturally viable without a pretrained language model.
\end{abstract}

\section{Introduction}

A frontier language model is a function: weights in, logits out. It has
no state between calls, no body to sense, no memory to consolidate, no
self to model, and no developmental trajectory. These are not missing
features in the sense of ``not yet shipped'' --- they are absent from
the bare architecture, and can only be approximated by external
scaffolding: state stores, retrieval indexes, agent loops, robot bodies.

Genesis was built to demonstrate that this is one architectural choice,
not a necessity of cognition. It is a two-layer cognitive architecture
in which cognition-relevant state is \emph{native}: a continuous
neurochemical substrate that evolves tick to tick, a semantic concept
network that grows through experience, a generative self-model that
predicts the system's own trajectory, and a closed control loop with
the physical hardware the system runs on. The system learns
continuously --- from conversation, teaching, autonomous reading, and
voluntary practice --- and is never retrained.

This paper makes the following contributions:

\begin{itemize}
  \item A description of a complete, running cognitive architecture
        that composes language and reasoning from a self-built semantic
        graph with no large language model anywhere in the system.
  \item A coupled neuromodulatory substrate (eighteen chemicals,
        nonlinear 18$\times$18 coupling, receptor desensitization and
        internalization, HPA axis, circadian oscillator) from which
        mental states --- flow, stress, sleep, overwhelm --- emerge
        rather than being selected.
  \item A machine-native embodiment loop: hardware interoception mapped
        to neurochemistry by functional equivalence, and neurochemical
        state mapped back to hardware control (CPU frequency,
        scheduling, I/O priority), including self-regulation of the
        cognitive process's own scheduling via emergent brain-wave
        state.
  \item An autopoietic active-inference loop: a generative model of the
        system's own neurochemical dynamics whose precision-weighted
        prediction errors feed back into those dynamics.
\end{itemize}

We are explicit about scope: Genesis is not a chatbot, not a
simulation of embodiment, and not a product. Its
language is narrower and less fluent than a trained language model's.
The contribution is architectural: the properties above follow from the
design, not from scaffolding around a frozen function.

\section{System Architecture}

\subsection{Two layers, one contract}

The system comprises a Rust daemon (the subcognitive layer) and a
Python process (the cognitive mind), communicating through a versioned
Unix-socket IPC protocol and a shared memory-mapped state file.

The Rust daemon is reactive (mind-driven dispatch) and owns all
dynamics that must run continuously: neurochemical integration,
receptor adaptation, HPA cascade, circadian phase, short-term to
long-term memory consolidation, association and dream synthesis, the
active-inference self-model, dyadic affective coupling, hardware
interoception, and body-control recommendations.

The Python layer owns cognition: perception, deliberation, a
compositional generative grammar, a concept network with typed
relations, graph-traversal reasoning, global workspace broadcast,
executive control, metacognition, theory of mind, narrative identity,
autonomous learning, and volitional drives (art, puzzles).

The boundary principle is strict: the substrate defines \emph{what
state exists and how it changes}; the cognitive mind defines \emph{how
that state is used}. Any cognitive module can be replaced --- e.g., a
rule-based perception classifier for a learned model --- without
touching the substrate, provided it reads the same neurochemical state
and writes to the same concept network.

\subsection{The core state: 3{,}288 bytes}

The entire runtime state is a single fixed-size \texttt{repr(C)} struct,
memory-mapped and shared across processes (Table~\ref{tab:layout}).
There are no heap allocations and no variable-length fields; readers
use a seqlock protocol and never see torn writes; a CRC32 checksum
detects corruption. All field offsets are asserted at compile time.

\begin{table}[h]
\centering
\begin{tabular}{@{}rlrl@{}}
\toprule
Offset & Field & Size & Contents \\
\midrule
0    & header              & 56    & magic, version, seqlock, identity \\
56   & neurochemicals      & 2{,}416 & 18 chemicals + coupling matrix \\
2{,}472 & zones            & 96    & task zone, emergent phase \\
2{,}568 & memory           & 120   & store pointers, gating weights \\
2{,}688 & manifest         & 536   & runtime module table \\
3{,}224 & checksum         & 4     & CRC32 \\
3{,}228 & inference signals& 60    & active-inference projection \\
\midrule
     & \textbf{total}      & \textbf{3{,}288} & (allocated: one 4\,KB page) \\
\bottomrule
\end{tabular}
\caption{Binary layout of the memory-mapped core state.}
\label{tab:layout}
\end{table}

This is not a configuration object; it is the brain state, addressable
to the byte, that every module reads and writes.

\section{The Neuromodulatory Substrate}

\subsection{Coupled dynamics}

Eighteen neurochemicals are modeled, each carrying level, baseline,
receptor sensitivity, velocity, receptor subtypes, desensitization and
internalization state, tonic/phasic components, and a vesicular pool.
A plastic 18$\times$18 coupling matrix defines mutual modulation with
three nonlinear mechanisms:

\begin{enumerate}
  \item \textbf{Hill-function threshold gating} (cooperativity $n{=}2$):
        coupling follows a sigmoid dose-response curve rather than
        linear scaling, modeling EC50 thresholds and receptor
        cooperativity.
  \item \textbf{State-dependent gain modulation}: coupling is amplified
        up to $+30\%$ when two chemicals co-activate and attenuated up
        to $-30\%$ under antagonism.
  \item \textbf{Subtype-aware dopamine coupling}, distinguishing
        dopaminergic pathways with different functional roles.
\end{enumerate}

Canonical interactions are reproduced structurally: cortisol suppresses
BDNF (chronic stress blocks plasticity), GABA inhibits glutamate,
endorphins disinhibit dopamine, and serotonin modulates BDNF recovery
(the mechanism behind delayed antidepressant onset). Emergent behavior
--- phase transitions among flow, stress, sleep, and overwhelm ---
arises from the coupling, not from any single variable.

\subsection{Receptor adaptation}

Effective level is not concentration:
\[
\mathrm{effective} = \mathrm{level} \times \mathrm{sensitivity}
\times f_{\mathrm{desens}} \times f_{\mathrm{intern}} .
\]
Fast desensitization (minutes) and slow internalization (hours) operate
independently, with floor values calibrated to published maxima:
receptor density floor 0.5, desensitization floor 0.2, internalization
floor 0.45. A system can therefore be functionally deficient in a
chemical it has in abundance --- tolerance, applied to neurochemistry.
Receptors resensitize during sleep.

\subsection{Metaplasticity}

The coupling matrix itself is plastic and lives in the shared state:
sustained co-activation patterns reshape the matrix, which reshapes
future dynamics --- the system learns how it learns. This provides a
structural basis for emotional development rather than a parameter
schedule.

\subsection{HPA axis and circadian drive}

Cortisol and CRH are modeled as stress hormones (zero at rest, released
on demand, decaying without a stressor), not homeostatic
neurotransmitters. The CRH$\to$ACTH$\to$cortisol cascade is gated by a
maturation level modeling the stress hyporesponsive period, preventing
neurotoxic cortisol exposure during development. A circadian oscillator
(SCN analogue) modulates baselines over a 24-hour cycle. Adenosine
accumulates with wakefulness and clears during sleep
\cite{borbely1982,xie2013}, so sleep pressure is emergent rather than
scheduled.

\section{Autopoietic Active Inference}

Before each neurochemical update, a generative model predicts the next
eighteen effective levels:
\[
\hat{x}_{t+1}[i] = \sum_j A[i][j]\,x_t[j] + b[i],
\]
with $A$ an 18$\times$18 transition matrix (identity-initialized) and
$b$ an offset vector (zero-initialized), updated by a precision-weighted
delta rule. Post-update prediction error (surprise) is fed back into
the dynamics it models --- the model is the system it predicts
\cite{hofstadter2007,friston2010}.

Precision adapts: sustained low surprise raises confidence; sustained
high surprise lowers it. Free energy combines surprise with a
precision-uncertainty term; expected future free energy drives
\emph{allostatic load}, which upregulates the cortisol baseline ---
anticipatory stress. High expected free energy also emits small
impulses toward predicted homeostatic states (active inference, not
merely prediction \cite{friston2009active,sterling2012}), and drives
neuromodulatory feedback: positive prediction error produces dopamine
(reward prediction error \cite{schultz1997,schultz2016}), high
surprise produces norepinephrine (orienting \cite{astonjones2005}) and a
metaplasticity
boost. The model persists across restarts and matures over time.

\section{Machine-Native Embodiment}

The embodiment is control-theoretic and machine-native. Hardware
signals map to neurochemistry by
\emph{functional equivalence}: a signal maps to a cognitive/affective
process when it has the same operational meaning for the machine's
survival and performance.

\paragraph{Afferent direction (interoception).} CPU temperature maps to
thermal state; CPU frequency to arousal/processing rate; own-process
CPU, RSS, and I/O to effort, cognitive load, and data-exchange activity;
battery to energy reserve; fan speed to thermoregulatory effort; power
draw to metabolic rate; RAPL power domains to spatial switching
activity (core$\to$NE, DRAM$\to$ACh); and perf counters
(cache/branch-miss ratios on the system's own process tree) to
microarchitectural prediction errors. Body-state signals (temperature,
battery) are global; activity signals are measured from the system's
own process tree only, so unrelated programs cannot induce chronic
allostatic load. Per-module telemetry is self-reported by a 25\,Hz
stack sampler that credits whichever subsystem's code is executing ---
a software EEG over the manifest's brain parts.

\paragraph{Efferent direction (body control).} The daemon publishes
recommendations: dopamine raises CPU frequency, GABA lowers it,
melatonin selects powersave; acetylcholine sets scheduling priority;
the plasticity gate sets I/O class; temperature imposes a hard thermal
cap regardless of dopaminergic drive. Application is cognitive-mind
driven: the cognitive process reads the recommendation, blends it with
its emergent brain-wave state (70\% cortical, 30\% body), and applies
its own scheduling decision --- cortical control of resource
allocation. Gamma dominance can override a body recommendation to
deprioritize, matching central autonomic network integration.

The loop is closed and physical: the system speeds itself up, which
raises its temperature, which caps its frequency, which slows it down.
It measures the effects of its own adjustments.

\section{Memory, Sleep, and Bounded Growth}

Short-term memory is a fixed-size ring buffer; consolidation promotes
high-salience entries to long-term storage gated by BDNF (stress
impairs consolidation, verified by integration test). Associative
retrieval uses SimHash matching; the association engine also drives
dream synthesis during sleep.

Memory grows during wakefulness and is bounded by three mechanisms run
during N3 sleep: (i) a vector-quantization codebook ($K{=}2048$
prototypes + int8 residuals replacing float32 embeddings), (ii) a
holographic associative graph (HRR, $D{=}2048$, 16 buckets per
relation; fixed $\sim$1.9\,MB regardless of edge count), and (iii)
sleep compression --- quantization, holographization of bridge edges,
LTM compaction via replay-and-bind, and global synaptic downscaling
\cite{plate1995,tononi2003}, executed as an atomic pass.

Sleep itself is emergent from adenosine accumulation ($\sim$16\,h to
threshold, $\sim$8\,h to clear, matching the circadian period). During
sleep a thalamic gate closes to background noise while salient stimuli
(direct address) can still wake the system; receptors resensitize;
consolidation and dream synthesis run.

\section{The Cognitive Layer}

The Python layer implements cognition on top of the substrate. Modules
grounded in cognitive theory include: a concept network (typed semantic
graph with STDP and spreading activation); graph-traversal reasoning
(transitive chains, causal inference, contradiction detection);
Dehaene-style global workspace broadcast with a 3--4 item bottleneck
\cite{dehaene2011,cowan2001}; a Damasio three-tier self (proto-self, core self,
autobiographical self \cite{damasio1999}); drift-diffusion decision
making \cite{ratcliff2008}; executive control; a recursive
metacognitive model that predicts its own reflection outcomes and
spawns/prunes levels by predictive usefulness; theory of mind; and a
narrative self-model with emergent personality traits and
Erikson-style developmental staging \cite{erikson1968}.

Language is compositional and generative: responses are assembled from
grammar structures weighted by emotional state and voice, over content
drawn from the concept network. A hard architectural invariant forbids
hardcoded responses --- seeds (vocabulary, grammar, relation verbs) are
permitted as building blocks, but the final utterance must be composed,
never recited. Dyadic coupling models the user's affective state as
hidden variables inferred from a VADER-style sentiment analyzer \cite{vader2014}, with
oxytocin-mediated attunement and rolling valence synchrony
\cite{palumbo2017}; the system
is tuned to a single primary user as the foundation of its
developmental model.

\section{Verification}

Over two thousand tests check the architecture's claims rather than
incidental code: binary layout offsets, seqlock consistency, coupled
dynamics, receptor adaptation, emergent phase transitions,
stress-impaired memory, HPA recovery, sleep emergence from adenosine,
burnout$\to$sleep$\to$recovery integration, daemon lifecycle, IPC
protocol, active-inference convergence, and property tests (levels in
$[0,1]$, no NaN after 1000 ticks, bounded coupling).

\section{Related Work}

\textbf{Cognitive architectures.} SOAR \cite{laird2012}, ACT-R
\cite{anderson1998}, and CLARION \cite{sun2002} established
production-system and hybrid symbolic-subsymbolic cognition; LIDA
\cite{franklin2016} implemented global workspace theory. Genesis differs in grounding all
modulation in a neuromodulatory substrate coupled to physical hardware,
and in its autopoietic self-model.

\textbf{Active inference.} The free energy principle
\cite{friston2010,friston2009active} and
predictive processing motivate the self-model; Genesis applies
precision-weighted prediction error to the system's own interoceptive
state rather than to external sensory streams --- a ``strange loop''
instantiation.

\textbf{Neuromodulation in AI.} Computational psychiatry and
biophysical models establish the receptor and coupling mechanisms
individually; we are not aware of a mainstream AI system integrating
them into a closed-loop substrate coupled to hardware control.

\textbf{Autonomous machine intelligence.} LeCun's proposed architecture
\cite{lecun2022}
(configurator, world model, cost, actor) argues for intrinsic-state
modulation of cognition; Genesis is a running instantiation of a
related position, with neuromodulation playing the role of the
intrinsic cost module.

\section{Limitations}

Language coverage is narrow relative to trained models; several
cognitive-layer modules (perception, question composition, memory
retrieval) are rule-based placeholders designed for replacement under
the fixed IPC contract. The neurochemical model is a neuromodulatory
abstraction, not a neural simulation --- no spiking activity, no
circuit anatomy. The dyadic model assumes a single primary user. The
system is one instance on one machine; no fleet-scale or multi-agent
claims are made.

\section{Conclusion}

Genesis demonstrates that continuous learning, emergent mental states,
bounded memory, hardware embodiment, self-modeling, and generative
language can coexist in a single local architecture with no pretrained
model --- because these properties follow from the design rather than
being scaffolded onto a frozen function. The system is running, the
verification suite and its coverage are stated, and the gaps are
named. The question the project poses is not whether this approach
matches frontier language models at their own task, but whether the
field has correctly identified which properties of minds are
architectural necessities and which are accidents of one successful
design.

\begin{thebibliography}{24}

\bibitem{friston2010}
K.~Friston.
\newblock The free-energy principle: a unified brain theory?
\newblock \emph{Nature Reviews Neuroscience}, 11:127--138, 2010.

\bibitem{friston2009active}
K.~Friston, J.~Daunizeau, and S.~J. Kiebel.
\newblock Reinforcement learning or active inference?
\newblock \emph{PLoS ONE}, 4(7):e6421, 2009.

\bibitem{schultz2016}
W.~Schultz.
\newblock Dopamine reward prediction error coding.
\newblock \emph{Dialogues in Clinical Neuroscience}, 18(1):23--32, 2016.

\bibitem{schultz1997}
W.~Schultz, P.~Dayan, and P.~R. Montague.
\newblock A neural substrate of prediction and reward.
\newblock \emph{Science}, 275(5306):1593--1599, 1997.

\bibitem{astonjones2005}
G.~Aston-Jones and J.~D. Cohen.
\newblock An integrative theory of locus coeruleus-norepinephrine
  function: adaptive gain and optimal performance.
\newblock \emph{Annual Review of Neuroscience}, 28:403--450, 2005.

\bibitem{sterling2012}
P.~Sterling.
\newblock Allostasis: a model of predictive regulation.
\newblock \emph{Physiology \& Behavior}, 106(1):5--15, 2012.

\bibitem{dehaene2011}
S.~Dehaene and J.-P. Changeux.
\newblock Experimental and theoretical approaches to conscious
  processing.
\newblock \emph{Neuron}, 70(2):200--227, 2011.

\bibitem{damasio1999}
A.~Damasio.
\newblock \emph{The Feeling of What Happens: Body and Emotion in the
  Making of Consciousness}.
\newblock Harcourt Brace, 1999.

\bibitem{palumbo2017}
R.~V. Palumbo, M.~E. Marraccini, L.~L. Weyandt, O.~Wilder-Smith,
  H.~A. McGee, S.~Liu, and M.~S. Goodwin.
\newblock Interpersonal autonomic physiology: a systematic review of
  the literature.
\newblock \emph{Personality and Social Psychology Review},
  21(2):99--141, 2017.

\bibitem{vader2014}
C.~J. Hutto and E.~Gilbert.
\newblock VADER: a parsimonious rule-based model for sentiment analysis
  of social media text.
\newblock In \emph{Proceedings of the Eighth International AAAI
  Conference on Weblogs and Social Media (ICWSM-14)}, pp.~216--225,
  2014.

\bibitem{plate1995}
T.~A. Plate.
\newblock Holographic reduced representations.
\newblock \emph{IEEE Transactions on Neural Networks}, 6(3):623--641,
  1995.

\bibitem{lecun2022}
Y.~LeCun.
\newblock A path towards autonomous machine intelligence.
\newblock OpenReview preprint, version 0.9.2, 2022.
  \url{https://openreview.net/forum?id=BZ5a1r-kVsf}

\bibitem{laird2012}
J.~E. Laird.
\newblock \emph{The Soar Cognitive Architecture}.
\newblock MIT Press, 2012.

\bibitem{anderson1998}
J.~R. Anderson and C.~Lebiere.
\newblock \emph{The Atomic Components of Thought}.
\newblock Lawrence Erlbaum Associates, 1998.

\bibitem{sun2002}
R.~Sun.
\newblock \emph{Duality of the Mind: A Bottom-Up Approach Toward
  Cognition}.
\newblock Lawrence Erlbaum Associates, 2002.

\bibitem{franklin2016}
S.~Franklin, T.~Madl, S.~Strain, U.~Faghihi, D.~Dong, S.~Kugele,
  J.~Snaider, P.~Agrawal, and S.~Chen.
\newblock A LIDA cognitive model tutorial.
\newblock \emph{Biologically Inspired Cognitive Architectures},
  16:105--130, 2016.

\bibitem{cowan2001}
N.~Cowan.
\newblock The magical number 4 in short-term memory: a reconsideration
  of mental storage capacity.
\newblock \emph{Behavioral and Brain Sciences}, 24(1):87--114, 2001.

\bibitem{hofstadter2007}
D.~R. Hofstadter.
\newblock \emph{I Am a Strange Loop}.
\newblock Basic Books, 2007.

\bibitem{borbely1982}
A.~A. Borb\'ely.
\newblock A two process model of sleep regulation.
\newblock \emph{Human Neurobiology}, 1(3):195--204, 1982.

\bibitem{tononi2003}
G.~Tononi and C.~Cirelli.
\newblock Sleep and synaptic homeostasis: a hypothesis.
\newblock \emph{Brain Research Bulletin}, 62(2):143--150, 2003.

\bibitem{xie2013}
L.~Xie, H.~Kang, Q.~Xu, M.~J. Chen, Y.~Liao, M.~Thiyagarajan,
  J.~O'Donnell, D.~J. Christensen, C.~Nicholson, J.~J. Iliff,
  T.~Takano, R.~Deane, and M.~Nedergaard.
\newblock Sleep drives metabolite clearance from the adult brain.
\newblock \emph{Science}, 342(6156):373--377, 2013.

\bibitem{ratcliff2008}
R.~Ratcliff and G.~McKoon.
\newblock The diffusion decision model: theory and data for two-choice
  decision tasks.
\newblock \emph{Neural Computation}, 20(4):873--922, 2008.

\bibitem{erikson1968}
E.~H. Erikson.
\newblock \emph{Identity: Youth and Crisis}.
\newblock Norton, 1968.

\end{thebibliography}

\end{document}
